\chapter{State of the Art}
\label{ch:sota}

\section{Neural Audio Synthesis}
\label{sec:sota_neural_audio}

The synthesis of audio using neural networks has evolved substantially over the past decade, transitioning from rule-based signal processing approaches toward data-driven models capable of producing perceptually convincing musical instrument sounds.

\paragraph{Early deep learning approaches.}
One of the foundational contributions in this area is NSynth \cite{engel2017nsynth}, which introduced the first large-scale neural audio synthesis system. NSynth uses a WaveNet autoencoder trained on a dataset of over 300{,}000 isolated musical notes across a wide range of instruments and pitches. The model learns a 16-dimensional embedding for each note and can interpolate between instrument timbres. However, NSynth is not designed for real-time operation: inference requires several minutes of compute time per second of audio, making it impractical for live performance applications.

\paragraph{Differentiable signal processing.}
A more interpretable and computationally efficient approach was introduced with DDSP \cite{engel2020ddsp}. Rather than generating raw audio samples or spectral frames, DDSP combines a neural network with explicit digital signal processing components~--- specifically, a harmonic synthesizer and a filtered noise model. The network predicts the control parameters (fundamental frequency, loudness, noise magnitudes) and the synthesizer generates audio from them. Because the synthesis model respects physical constraints of sound production, DDSP requires far less data and computation than purely data-driven approaches. It has been applied successfully to violin and flute synthesis, and later to voice-to-instrument conversion. The key limitation for our use case is that DDSP requires explicit pitch annotations during training (or a pitch tracker at inference time), whereas RAVE-based approaches learn pitch implicitly from raw audio. DDSP also lacks a dedicated real-time streaming architecture for Pure Data deployment.

\paragraph{RAVE: real-time variational autoencoder.}
The most directly relevant prior work is RAVE \cite{caillon2021rave}. RAVE is a two-stage generative model based on a Variational Autoencoder (VAE). In the first stage, an encoder compresses audio into a compact latent representation and a decoder reconstructs audio from it. The encoder uses multi-band decomposition for efficiency, and the decoder operates in the spectral domain via a multi-scale STFT reconstruction loss rather than generating raw samples. In the second stage, a normalizing flow prior is trained over the latent space to enable unconditional generation. The key contribution of RAVE is that the encoder-decoder pair is compact enough to run in real-time: the authors demonstrate inference at 40$\times$ real-time speed on a consumer CPU, enabling live performance applications. RAVE is deployable in Pure Data via the \texttt{nn\textasciitilde} external, making it the standard choice for neural timbre transfer in live contexts.

\paragraph{BRAVE: streaming real-time RAVE.}
BRAVE \cite{caspe2025brave} extends RAVE with a \emph{streaming encoder}~--- a causal architecture that processes audio in small temporal chunks rather than fixed-size windows. This architectural change reduces end-to-end latency significantly while preserving audio quality, making BRAVE better suited than RAVE for real-time performance scenarios where latency is perceptible. The authors evaluate BRAVE on two datasets: a solo saxophone corpus (Filosax, $\sim$4.5 hours) and an electronic drum kit corpus ($\sim$2.8 hours). Results show that BRAVE achieves competitive perceptual quality with RAVE while providing lower latency. An important limitation noted in the paper is the difficulty of accurate pitch rendering for melodic instruments: while the model preserves relative pitch contours, absolute pitch accuracy degrades at extreme registers. This limitation is particularly relevant to our voice-to-guitar application and constitutes one of the research questions we investigate.

\section{Timbre Transfer and Voice Conversion}
\label{sec:sota_timbre}

\paragraph{What is timbre?}
Timbre refers to the perceptual quality that distinguishes two sounds with the same pitch and loudness~--- what makes a guitar sound different from a violin playing the same note. Acoustically, timbre is determined by the spectral envelope (distribution of energy across frequencies), the attack and decay characteristics of notes, and fine-grained temporal modulations. Neural timbre transfer aims to learn a mapping from the timbre of one instrument to another while preserving other musical properties such as pitch and rhythm.

\paragraph{CycleGAN-based approaches.}
TimbreTron \cite{huang2018timbretron} applied CycleGAN~--- a generative adversarial network designed for unpaired image-to-image translation~--- to the problem of timbre transfer. Operating on constant-Q transform (CQT) spectrograms, TimbreTron learns to convert between instrument timbres without requiring paired recordings. While the approach demonstrates promising results for piano-to-harpsichord and classical-to-jazz stylistic transfers, it suffers from characteristic GAN instabilities, requires significant data, and does not operate in real-time. The inference involves spectrogram-domain manipulation followed by WaveNet-based vocoding, making it unsuitable for live applications.

\paragraph{Voice conversion.}
Voice conversion (VC) is the closely related task of modifying a speaker's voice to sound like a different speaker while preserving linguistic content. Modern VC systems such as AutoVC \cite{qian2019autovc} and SoundStream-based approaches use encoder-decoder architectures with speaker conditioning, achieving high-quality results in offline settings. The cross-modal variant~--- converting voice to a non-voice instrument~--- is less studied. The primary challenge is that the acoustic feature spaces of voice and instrument differ in fundamental ways: voice has formant structure driven by the vocal tract, while instruments have body resonance, mechanical excitation, and no phonetic content. RAVE-based approaches sidestep this mismatch by learning a shared latent space that captures timbral structure without explicitly modelling either source.

\section{Datasets for Neural Audio Synthesis}
\label{sec:sota_datasets}

The quality and quantity of training data are critical factors in RAVE-based synthesis. The model requires continuous, single-instrument audio without silence or background noise, and benefits from diverse musical content that covers a wide range of pitches, dynamics, and playing techniques.

\paragraph{GuitarSet}
\cite{xi2018guitarset} is a dataset of 360 recordings of acoustic guitar performance by six players across five musical styles. Each recording ($\sim$30~seconds) is captured with both a hexaphonic pickup and a reference microphone, providing rich annotations including pitch contours, beat positions, chord labels, and playing technique markers. The dataset is released under CC~BY~4.0.

\paragraph{Guitar-TECHS}
is a dataset of electric guitar performances recorded in multiple environments by three professional guitarists. It covers six playing techniques, musical excerpts, chord types, and major scale runs, providing multiple simultaneous capture channels including direct-injection (DI), amplified microphone, and binaural room recordings. Released under CC~BY~4.0.

\paragraph{IDMT-SMT-Guitar}
\cite{kehling2014idmtsmtguitar} provides large-scale guitar recordings covering single-note classification (datasets~1 and~2), short licks (dataset~3), and extended performance excerpts across various playing styles and tempos (dataset~4). For continuous-audio training, datasets~4 and~3 are the most useful subsets. The dataset includes both acoustic and electric guitars, recorded through multiple capture paths (microphone, pickup, direct input).

\paragraph{Groove MIDI Dataset}
\cite{gillick2019groove} is a corpus of human drum performances by professional drummers, captured as MIDI and audio. The audio portion consists of approximately 13 hours of drum kit recordings at 44.1~kHz, covering a wide range of musical styles and tempos. Released under CC~BY~4.0.

\section{Research Gap and Motivation}
\label{sec:sota_gap}

The review above reveals a clear gap: while BRAVE has been applied to saxophone and drums, its application to guitar timbre transfer from voice has not been systematically investigated. The specific challenge of voice-to-guitar conversion raises open questions that the BRAVE authors acknowledge but do not address:

\begin{enumerate}
    \item \textbf{Pitch rendering for plucked string instruments.} Guitar sound is characterised by a sharp percussive attack followed by exponential decay. Unlike saxophone, where breath pressure sustains a note, guitar notes cannot be sustained indefinitely. How this temporal envelope interacts with the streaming encoder's representation is unknown.
    \item \textbf{Timbral gap between voice and guitar.} The acoustic feature space of singing voice differs substantially from guitar. Voice contains formant structure, vibrato, and phonetic content absent from guitar. It is unclear how much of this mismatch the BRAVE latent space can absorb.
    \item \textbf{Multi-source training.} Training BRAVE on a combined acoustic plus electric guitar corpus with heterogeneous recording paths has not been explored. Whether this improves generalisation or degrades timbre fidelity is an empirical question.
\end{enumerate}

This project addresses these gaps by training and evaluating BRAVE on the guitar corpus described above, documenting the qualitative and quantitative characteristics of the resulting timbre transfer, and building a real-time Pure Data system that demonstrates the approach in a live performance context. The secondary beatbox-to-drums experiment, replicating the drums evaluation from the original BRAVE paper on our setup, provides a methodological baseline.
